% BLOCK14_STAGE2_MANAGED
% Lecture 14: STL containers, iterators, algorithms, complexity and invalidation
\section{Лекция 14. STL: контейнеры, итераторы и алгоритмы}
\label{sec:lecture-14}

До этого мы изучали отдельные части standard library:
\texttt{std::vector}, \texttt{std::map}, smart pointers, \texttt{optional},
lambdas и templates. Теперь объединим их в главную архитектурную идею STL:
\[
\text{container}
\rightarrow
\text{iterator/range}
\rightarrow
\text{algorithm}
\rightarrow
\text{callable policy}.
\]

В прикладной части будем строить pipeline обработки каталога учебных блоков:
хранение записей, индекс, поиск, фильтрация, сортировка, преобразование,
агрегация и удаление. Одновременно разберём complexity guarantees и invalidation,
потому что выбор container определяется не только удобством синтаксиса.

\begin{importantbox}[Идея STL]
Container отвечает за хранение и lifetime элементов. Iterator описывает позицию
и способ пройти range. Algorithm работает с range, не обязуясь знать конкретный
container. Lambda/function object задаёт переменную часть поведения.
\end{importantbox}

\subsection{Что такое STL в этом блоке}

Исторически термин STL относился к Standard Template Library, а современная
C++ standard library шире. В учебной практике под «STL-стилем» мы будем иметь в
виду прежде всего:
containers; iterators; algorithms; function objects/predicates; связанные utility abstractions.

Не вся standard library относится к STL, и не нужно смешивать эти понятия как
строгие синонимы.

\subsection{Почему контейнеров несколько}

Если бы единственным критерием было «хранить набор объектов», одного
\texttt{vector} хватило бы почти всегда. Но программы предъявляют разные
требования:
нужен ли contiguous storage; важен ли порядок; нужен ли быстрый lookup по key; часто ли вставляем в середину; нужны ли стабильные references/iterators; сколько allocations приемлемо; важна ли cache locality; нужен ли deterministic ordered traversal.

Выбор container --- часть design и cost model.

\subsection{\texttt{std::array}}

\texttt{std::array<T, N>} --- fixed-size container, размер известен в type.

\begin{cppcode}
std::array<int, 4> values{1, 2, 3, 4};
\end{cppcode}

Плюсы:
contiguous elements; size не меняется; обычный container API: \texttt{begin()}, \texttt{end()}, \texttt{size()}; хорошо работает с algorithms.

Это удобная modern abstraction над built-in array, когда размер фиксирован.

\subsection{\texttt{std::vector}}

\texttt{vector<T>} --- основной general-purpose sequence container:
элементы contiguous; быстрый random access; amortized constant-time \texttt{push\_back}; обычно хорошая locality при sequential traversal; рост может вызвать reallocation.

\begin{cppcode}
std::vector<CourseBlock> blocks;
blocks.push_back({14, "stl", 2});
\end{cppcode}

Для большинства динамических последовательностей начинайте с \texttt{vector}, если нет конкретного требования, которое ему противоречит.

\subsection{Size и capacity}

У \texttt{vector} различаются:
\texttt{size()} --- число живых элементов; \texttt{capacity()} --- сколько элементов можно разместить без следующего расширения storage.

\begin{cppcode}
std::vector<int> values;
values.reserve(100);
\end{cppcode}

\texttt{reserve(100)} не создаёт 100 элементов и не меняет \texttt{size()}.
Он просит capacity не меньше 100.

\begin{standardbox}[Что гарантирует reserve]
Если requested capacity больше текущей, происходит reallocation. Если нет ---
reallocation не происходит. После успешного \texttt{reserve} вставки не
вызывают reallocation, пока size не превысит capacity.
\end{standardbox}

\subsection{Почему reallocation важен}

При росте \texttt{vector} может получить новый storage и перенести/скопировать
элементы. Старый storage освобождается.

Следствия:
addresses элементов меняются; pointers/references/iterators на элементы становятся invalid при reallocation; стоимость отдельного роста может быть линейной по числу элементов; последовательность многих \texttt{push\_back} всё равно имеет amortized constant complexity.

\begin{warningbox}[Amortized O(1) не означает каждый push\_back O(1)]
Иногда \texttt{push\_back} требует allocation и переноса существующих элементов.
Гарантия говорит о средней стоимости последовательности операций в
амортизированном анализе.
\end{warningbox}

\subsection{\texttt{std::deque}}

\texttt{deque<T>} --- sequence с efficient insertion/removal на обоих концах.

Он не обязан хранить все элементы в одном contiguous block. Random access
поддерживается, но locality/layout отличаются от \texttt{vector}.

Обычный выбор:
часто нужны \texttt{push\_front}/\texttt{pop\_front}; random access всё ещё нужен; contiguous storage не требуется.

Не делайте вывод «deque = vector без reallocation»: правила invalidation и
representation другие и должны читаться отдельно.

\subsection{\texttt{std::list}}

\texttt{list<T>} --- doubly linked list.

Полезные свойства:
insertion/erase по известному iterator обычно constant complexity; references/iterators к другим элементам обычно сохраняются при insertion/erase; нет random-access iterator; каждый node обычно хранится отдельно; locality при traversal обычно хуже contiguous vector.

\begin{performancebox}[O(1) insertion не делает list автоматически быстрее]
Чтобы вставить «в середину по номеру», сначала может понадобиться линейно найти
позицию. Отдельные allocations и pointer chasing также дают большие constant
factors. Asymptotics --- только часть стоимости.
\end{performancebox}

\subsection{\texttt{std::forward\_list}}

\texttt{forward\_list<T>} --- singly linked list. Он экономнее по links и
поддерживает операции around a known predecessor, но interface менее удобен.

В обычном application code встречается реже. Выбирать его стоит под конкретное
требование, а не потому что «односвязный список проще».

\subsection{Ordered associative containers}

\texttt{std::map<Key, T>} и \texttt{std::set<Key>} поддерживают ordering по
comparator.

Обычные guarantees:
lookup/insert/erase --- logarithmic complexity; traversal идёт в порядке comparator; insertion обычно не invalidates iterators/references к существующим элементам; erase invalidates только erased elements.

\texttt{map} хранит key-value pairs, \texttt{set} --- keys.

\subsection{Unordered associative containers}

\texttt{std::unordered\_map} и \texttt{std::unordered\_set} организованы вокруг
hashing.

Обычный contract:
average-case lookup/insert/erase --- constant complexity; worst case может быть linear; iteration order не sorted; rehash влияет на iterators; quality hash/equality важна.

\begin{warningbox}[O(1) у unordered\_map --- average, не абсолютное обещание]
Нельзя писать «unordered\_map всегда быстрее map». Учитываются размер данных,
hash cost, allocations, locality, ordering requirements и worst-case behavior.
\end{warningbox}

\subsection{Когда map, когда unordered\_map}

\texttt{map} подходит, если:
нужен sorted traversal; нужны ordered operations вроде \texttt{lower\_bound}; важен logarithmic worst-case contract.

\texttt{unordered\_map} подходит, если:
нужен lookup по exact key; порядок не важен; хороший hash доступен; average constant lookup соответствует задаче.

\subsection{Адаптеры контейнеров}

\texttt{stack}, \texttt{queue}, \texttt{priority\_queue} --- container adapters.
Они намеренно предоставляют более узкий interface поверх underlying container.

Например, \texttt{queue} выражает FIFO semantics. Если application должен
доставать произвольный элемент по iterator, возможно, \texttt{queue} выбран
неверно.

\subsection{Iterator: абстракция позиции}

Iterator --- объект, с помощью которого algorithm обращается к position в
range.

\begin{cppcode}
auto it = blocks.begin();

if (it != blocks.end()) {
    std::cout << it->slug;
}
\end{cppcode}

Iterator не обязан быть raw pointer, хотя iterator contiguous container может
быть pointer-like.

Не нужно утверждать, что iterator «это указатель внутри STL». Это abstraction
с определёнными operations и guarantees.

\subsection{Half-open range \texttt{[first, last)}}

Большинство classic algorithms принимает пару iterators:
\[
[first, last).
\]

\texttt{first} относится к range, \texttt{last} указывает на позицию
\emph{после} последнего элемента.

Преимущества:
empty range задаётся \texttt{first == last}; length естественно выражается distance; соседние ranges стыкуются без специальной обработки границы.

\subsection{Iterator categories}

Не все iterators поддерживают одинаковые операции.

Классическая иерархия требований:
input iterator; output iterator; forward iterator; bidirectional iterator; random-access iterator; в C++20 отдельно формализован contiguous iterator.

\texttt{vector} предоставляет random-access iterators.
\texttt{list} --- bidirectional iterators.

Поэтому algorithm \texttt{std::sort} требует random-access iterator и не
работает с \texttt{std::list::iterator}.

\subsection{Почему list имеет member \texttt{sort}}

\begin{cppcode}
std::list<int> values{4, 1, 3};
values.sort();
\end{cppcode}

List знает свою node-based structure и может сортировать её без random-access
iterator. Попытка передать \texttt{list.begin()/end()} в \texttt{std::sort}
проверяется в практике как real compile failure.

\subsection{Const iterator}

\texttt{const\_iterator} не разрешает менять элемент через iterator.

\begin{cppcode}
const std::vector<int> values{1, 2, 3};
auto it = values.begin(); // iterator to const element
\end{cppcode}

Также есть \texttt{cbegin()}/\texttt{cend()}, если хочется явно запросить
read-only iterator у non-const container.

\subsection{Range-based for}

\begin{cppcode}
for (const auto& block : blocks) {
    std::cout << block.slug;
}
\end{cppcode}

Это удобный syntax поверх begin/end pattern, когда алгоритм прост:
последовательно пройти range.

Но если задача --- «найти», «отсортировать», «преобразовать», «посчитать»,
standard algorithm часто выражает намерение лучше ручного loop.

\subsection{\texttt{std::find} и \texttt{std::find\_if}}

\begin{cppcode}
auto it = std::find_if(
    blocks.begin(),
    blocks.end(),
    [](const CourseBlock& block) {
        return block.number == 14;
    }
);
\end{cppcode}

Result --- iterator. Если ничего не найдено, он равен \texttt{end()}.

Это non-owning position в container, поэтому применяются lifetime/invalidation
rules container.

\subsection{\texttt{std::count\_if}}

\begin{cppcode}
const auto ready = std::count_if(
    blocks.begin(),
    blocks.end(),
    [](const CourseBlock& block) {
        return block.complete;
    }
);
\end{cppcode}

Algorithm принимает predicate из блока~\ref{sec:lecture-13}. Container при этом
может быть другим, если iterator requirements выполнены.

\subsection{\texttt{std::copy} и output iterators}

Algorithms могут писать через output iterator.

\begin{cppcode}
std::vector<CourseBlock> ready;

std::copy_if(
    blocks.begin(),
    blocks.end(),
    std::back_inserter(ready),
    [](const CourseBlock& block) {
        return block.complete;
    }
);
\end{cppcode}

\texttt{std::back\_inserter} создаёт adapter, который превращает assignment
algorithm-а в \texttt{push\_back} container.

\subsection{\texttt{std::transform}}

\begin{cppcode}
std::vector<std::string> slugs;

std::transform(
    blocks.begin(),
    blocks.end(),
    std::back_inserter(slugs),
    [](const CourseBlock& block) {
        return block.slug;
    }
);
\end{cppcode}

Mechanism traversal находится в algorithm, conversion policy --- в lambda.

\subsection{\texttt{std::accumulate}}

\texttt{std::accumulate} находится в \headerfile{numeric}:

\begin{cppcode}
const int examples = std::accumulate(
    blocks.begin(),
    blocks.end(),
    0,
    [](int total, const CourseBlock& block) {
        return total + block.examples;
    }
);
\end{cppcode}

Тип initial value важен: он влияет на accumulator type.

Например, initial \texttt{0} --- \texttt{int}; initial \texttt{0.0} ---
\texttt{double}.

\subsection{\texttt{std::sort}}

\begin{cppcode}
std::sort(
    blocks.begin(),
    blocks.end(),
    [](const CourseBlock& left, const CourseBlock& right) {
        return left.number < right.number;
    }
);
\end{cppcode}

Для \texttt{std::sort} complexity --- $O(N \log N)$ comparisons в standard
guarantee современной C++ library.

Comparator должен задавать корректное strict weak ordering.

\subsection{Ошибка comparator}

Плохо:

\begin{cppcode}
[](int a, int b) {
    return a <= b;
}
\end{cppcode}

Для равных значений comparator возвращает true в обе стороны и нарушает
требования ordering.

Правильно для ascending order:

\begin{cppcode}
[](int a, int b) {
    return a < b;
}
\end{cppcode}

Algorithm contracts включают требования не только к iterator, но и к callable.

\subsection{Stable sort}

\texttt{std::stable\_sort} сохраняет relative order equivalent elements.
Обычный \texttt{std::sort} этого не обещает.

Если сначала отсортировать по secondary key, а затем stable-sort по primary key,
можно сохранить нужный secondary ordering. Но зачастую яснее сразу написать
comparator по tuple ключей.

\subsection{\texttt{lower\_bound}}

Для sorted range:

\begin{cppcode}
auto it = std::lower_bound(
    values.begin(),
    values.end(),
    target
);
\end{cppcode}

Algorithm использует ordering и на random-access iterator выполняет logarithmic
число comparisons. Но стоимость iterator movement зависит от iterator category.

Для ordered associative container часто разумнее использовать member
\texttt{lower\_bound}, потому что container может использовать внутреннюю
tree structure.

\subsection{\texttt{all\_of}, \texttt{any\_of}, \texttt{none\_of}}

Эти algorithms напрямую выражают логические queries:

\begin{cppcode}
const bool all_ready = std::all_of(
    blocks.begin(),
    blocks.end(),
    [](const CourseBlock& block) {
        return block.complete;
    }
);
\end{cppcode}

Они могут завершить обход раньше, как только result уже определён.

\subsection{Алгоритм не знает container}

Это основная сила STL.

\begin{cppcode}
template <typename Iterator, typename Predicate>
std::size_t count_matching(
    Iterator first,
    Iterator last,
    Predicate predicate);
\end{cppcode}

Concrete algorithm может работать с \texttt{vector}, \texttt{list}, built-in
array или другим range, если iterators удовлетворяют требованиям.

Templates дают static polymorphism по iterator/callable types.

\subsection{Complexity --- часть API}

Для container/algorithm complexity guarantee нужно читать как contract.

Примеры:
\texttt{vector::operator[]} --- constant; \texttt{vector::push\_back} --- amortized constant; \texttt{map::find} --- logarithmic; \texttt{unordered\_map::find} --- average constant; \texttt{find\_if} --- linear в размере range в worst case; \texttt{sort} --- $O(N \log N)$ comparisons.

Но Big-O не описывает cache misses, allocation cost, branch behavior и object
size.

\subsection{Cache locality}

Contiguous \texttt{vector} часто хорошо подходит sequential traversal:
соседние элементы расположены рядом в storage.

Node-based \texttt{list/map} обычно требует переходов между nodes, которые
могут находиться в разных местах памяти.

Это обычная performance model, а не стандартное обещание конкретного cache
поведения CPU. Измерять нужно на реальном workload.

\subsection{Iterator invalidation}

\term{Invalidation} означает, что ранее полученный iterator/reference/pointer
больше нельзя использовать в тех operations, для которых он был валиден.

Причина --- container изменил storage или удалил соответствующий element.

Ошибки invalidation особенно опасны, потому что source code выглядит невинно:

\begin{cppcode}
auto it = values.begin();
values.push_back(42);
// *it может быть invalid, если случилась vector reallocation
\end{cppcode}

\subsection{Vector invalidation при росте}

Если \texttt{vector} reallocate:
invalid становятся все iterators; references; pointers на elements.

Если reallocation не было, insertion rules всё равно могут invalidировать
positions в зависимости от места операции. Например, insertion/erase в
середине сдвигают последующие elements.

Всегда читайте contract конкретной operation.

\subsection{Как безопасно наблюдать reallocation}

В практике мы сохраняем:
\texttt{data()} address как число/для сравнения до изменения; старый capacity; затем вызываем \texttt{push\_back}; после операции сравниваем старый и новый address; \emph{не разыменовываем} старый pointer.

Так можно показать факт reallocation без UB.

Отдельный специально небезопасный executable выполняется только под ASan и
должен дать diagnostic.

\subsection{List stability}

Для \texttt{list} insertion обычно не invalidates iterators/references к
существующим elements. Erase invalidates только iterator/reference на erased
element.

Это сильное отличие от \texttt{vector}, но за него платим node allocations,
отсутствием random access и locality.

\subsection{Map/set stability}

Для ordered associative containers insertion/emplacement не invalidates
iterators/references к существующим elements; erase invalidates erased
elements.

Unordered containers имеют дополнительную ось --- rehash. Rehash invalidates
iterators, хотя references/pointers к elements имеют отдельные guarantees.

Не переносите правила одного container на другой.

\subsection{Erase у vector}

\begin{cppcode}
auto it = values.erase(position);
\end{cppcode}

Result указывает на element, следующий за erased, либо \texttt{end()}.

После erase erased element уничтожен, а элементы после него сдвигаются.
Iterators/references на erased и subsequent elements invalidated.

\subsection{Erase-remove idiom в C++17}

Classic \texttt{std::remove\_if} сам не уменьшает size container.

\begin{cppcode}
auto new_end = std::remove_if(
    values.begin(),
    values.end(),
    [](int value) {
        return value < 0;
    }
);

values.erase(new_end, values.end());
\end{cppcode}

\texttt{remove\_if} переставляет retained elements к началу range и возвращает
logical new end. Реальное удаление хвоста делает container::erase.

C++20 добавляет удобные \texttt{std::erase}/\texttt{std::erase\_if} для
поддерживаемых containers, но основной стандарт блока --- C++17.

\subsection{Почему remove\_if не вызывает erase сам}

Generic algorithm знает только iterator range и не обязан иметь доступ к
container object или его size-management API.

Пример архитектуры STL:
algorithm работает с range, container отвечает за structural mutation.

\subsection{Алгоритм или member function}

Иногда container предоставляет member function с похожей задачей.

Примеры:
\texttt{list::sort} вместо \texttt{std::sort}; \texttt{map::find} использует tree structure; \texttt{unordered\_map::find} использует hash table.

Если member может использовать внутреннюю структуру container, он часто точнее
выражает задачу и даёт нужную complexity.

\subsection{Index рядом с vector}

Прикладной каталог удобно хранить в \texttt{vector<CourseBlock>} для
последовательной обработки, а дополнительный lookup index строить отдельно:

\begin{cppcode}
std::unordered_map<int, std::size_t> index;

for (std::size_t i = 0; i < blocks.size(); ++i) {
    index.emplace(blocks[i].number, i);
}
\end{cppcode}

Теперь:
vector остаётся source of truth; index хранит positions, а не copies objects; lookup по number может быть average O(1).

Но structural mutation vector может сделать сохранённые indices логически
неверными. Index нужно rebuild/update по contract.

\subsection{Почему index по position безопаснее pointer только условно}

Index \texttt{size\_t} не становится dangling address после reallocation.
Но если insert/erase меняет порядок элементов, старый index может указывать уже
на другой object.

Иными словами:
pointer invalidation --- memory/lifetime problem; stale index --- logic/invariant problem.

Оба случая требуют явного design.

\subsection{Pipeline: filter → sort → transform → aggregate}

Реальный processing pipeline можно писать по шагам:

\begin{cppcode}
std::vector<CourseBlock> selected;

std::copy_if(
    blocks.begin(),
    blocks.end(),
    std::back_inserter(selected),
    [](const CourseBlock& block) {
        return block.examples >= 2;
    }
);

std::sort(
    selected.begin(),
    selected.end(),
    [](const CourseBlock& a, const CourseBlock& b) {
        return a.number < b.number;
    }
);
\end{cppcode}

Дальше \texttt{transform} получает slugs, \texttt{accumulate} считает examples.

Промежуточный vector стоит памяти/copies, но делает pipeline прозрачным.
Не надо раньше времени превращать код в нечитаемую цепочку abstractions.

\subsection{Algorithms и mutation}

Некоторые algorithms только читают:
\texttt{find\_if}, \texttt{count\_if}, \texttt{all\_of}.

Другие меняют elements через iterators:
\texttt{sort}, \texttt{transform} in-place, \texttt{remove\_if}.

Но structural properties container, например size, generic algorithm обычно
не меняет сам.

Разделяйте element mutation и container structural mutation.

\subsection{Iterator lifetime и owner lifetime}

Iterator --- non-owning handle. Container должен продолжать существовать,
а operation не должна была invalidate iterator.

\begin{cppcode}
std::vector<int>::iterator bad();

std::vector<int>::iterator bad() {
    std::vector<int> local{1, 2, 3};
    return local.begin(); // iterator переживёт container
}
\end{cppcode}

Это тот же lifetime принцип, что raw pointer/reference из ранних блоков.

\subsection{End iterator тоже может invalidated}

Нельзя думать, что только iterators на «реальные элементы» важны.
Past-the-end iterator также относится к container state и может быть invalidated
операциями.

Поэтому сохранять \texttt{end()} до structural mutation и использовать после
неё без проверки contract опасно.

\subsection{Const-correct algorithms}

Если algorithm только читает, передавайте const range/container там, где это
естественно.

\begin{cppcode}
int total_examples(const std::vector<CourseBlock>& blocks) {
    return std::accumulate(
        blocks.cbegin(),
        blocks.cend(),
        0,
        [](int total, const CourseBlock& block) {
            return total + block.examples;
        }
    );
}
\end{cppcode}

Это делает невозможной случайную mutation элементов через iterator.

\subsection{Не копировать container без причины}

Параметр:

\begin{cppcode}
void print(std::vector<CourseBlock> blocks);
\end{cppcode}

создаёт своё value, что может означать copy/move container.

Если функция только читает:

\begin{cppcode}
void print(const std::vector<CourseBlock>& blocks);
\end{cppcode}

Но generic algorithm APIs часто лучше принимать iterator pair или template
range-like abstraction, если нужна независимость от конкретного container.

\subsection{Когда конкретный container type в API уместен}

Если функция использует vector-specific guarantees:
contiguous data; random access; capacity/reserve.

тогда \texttt{vector<T>\&} в API может быть честным contract.

Не нужно делать всё template-generic только ради generic programming.

\subsection{Exception safety}

Standard containers управляют своими resources через RAII. Но сильная/basic
exception guarantee конкретной operation зависит от type properties и самой
operation.

Например, vector reallocation может copy/move elements. \texttt{noexcept} move
из блока~\ref{sec:lecture-12} может влиять на выбранный library strategy.

Не утверждайте универсально «любая STL операция strong exception safe».

\subsection{Allocators: только первый взгляд}

Containers parameterized allocator type, например:

\begin{cppcode}
std::vector<int, Allocator>
\end{cppcode}

Allocator abstraction управляет получением/освобождением storage для container.
Полный allocator model сложен и будет кратко рассмотрен в блоке 15.

В обычном application code default allocator достаточен, пока profiling или
специальная memory architecture не доказали другое.

\subsection{Практический выбор container}

Начальная эвристика:
dynamic sequence --- \texttt{vector}; fixed sequence --- \texttt{array}; efficient both-end queue --- \texttt{deque}; node stability/splice-specific workload --- рассмотреть \texttt{list}; ordered key lookup --- \texttt{map}/\texttt{set}; unordered exact-key lookup --- \texttt{unordered\_map}/\texttt{unordered\_set}.

Потом проверяются:
complexity, invalidation, ordering, locality, allocations и actual workload.

\subsection{Эксперимент 1: course\_catalog\_pipeline}

Каталог:
\filepath{examples/14\_stl\_containers\_iterators\_algorithms/01\_course\_catalog\_pipeline}.

Программа строит прикладной каталог 15 blocks:
1) source of truth --- \texttt{vector<CourseBlock>}; 2) number index --- \texttt{unordered\_map<int, size\_t>}; 3) \texttt{find\_if} и index дают два способа query; 4) \texttt{copy\_if} фильтрует готовые blocks; 5) \texttt{sort} упорядочивает выборку; 6) \texttt{transform} получает slugs; 7) \texttt{accumulate} считает число examples; 8) \texttt{map} строит ordered summary по status; 9) \texttt{optional<CourseBlock>} представляет missing query.

\begin{shellcode}
cd examples/14_stl_containers_iterators_algorithms/01_course_catalog_pipeline
./run_checks.sh
\end{shellcode}

\subsection{Эксперимент 2: invalidation\_and\_erase\_lab}

Каталог:
\filepath{examples/14\_stl\_containers\_iterators\_algorithms/02\_invalidation\_and\_erase\_lab}.

Safe executable проверяет:
\texttt{vector::reserve} и growth без reallocation внутри capacity; forced reallocation и изменение \texttt{data()} address без использования старого pointer; erase-remove idiom; возвращаемый iterator \texttt{erase}; list iterator/reference stability при вставке другого element; unordered-map rehash observation без использования invalidated iterator.

Отдельный ASan-only executable намеренно разыменовывает pointer после forced
vector reallocation и должен дать use-after-free diagnostic.

Ещё один negative compile file подтверждает, что \texttt{std::sort} нельзя
применить к \texttt{std::list} iterators.

\begin{shellcode}
cd examples/14_stl_containers_iterators_algorithms/02_invalidation_and_erase_lab
./run_checks.sh
\end{shellcode}

\subsection{Что проверяет практика, а что нет}

Experiments подтверждают поведение текущей toolchain и конкретных inputs.
Они не заменяют standard guarantee.

Например:
изменение \texttt{data()} действительно показывает reallocation в case, который специально превышает capacity; ASan diagnostic показывает конкретную lifetime ошибку; отсутствие sanitizer report в safe run не доказывает отсутствие всех UB; measured addresses не превращаются в переносимое обещание allocator.

\subsection{Итоговый checklist STL}

Перед выбором container спросите:
какие операции доминируют; какой complexity contract нужен; нужен ли ordering; нужна ли pointer/reference stability; важна ли contiguous storage/locality; как часто меняется структура container.

Перед использованием iterator:
жив ли owner container; не было ли invalidating operation; удовлетворяет ли iterator category требованиям algorithm; не сохранили ли мы stale \texttt{end()}.

Перед написанием manual loop:
есть ли standard algorithm, который точнее выражает намерение; нужен ли predicate/lambda из блока 13; structural mutation отделена от element processing.

\begin{importantbox}[Главный вывод]
STL --- не список удобных функций. Это система contracts между storage,
positions, algorithms и callable policies. Хороший C++ code выбирает container
по требованиям, использует algorithms для намерения и всегда учитывает
complexity и invalidation.
\end{importantbox}

Следующий, последний базовый блок соберёт современный C++ в единую картину:
\texttt{constexpr}, C++20 Concepts, обзор coroutines и modules, allocators,
tooling, дальнейший маршрут обучения и критерии выбора возможностей языка.
